\section{Interactive Oracle Proof}

Now that we have covered the core construction behind our contribution, let's take a holistic view of verification of a pairing computation with an interactive proof system. While the construction above works beautifully on it's own, we can further optimize the protocol by batching all polynomial verifications into a two round interactive proof.

\subsection{Batching Polynomial Identity Testing}

The whole process of polynomial verification can be batched into a single polynomial verification. Here's how a two round interactive proof system would look like,

The prover can commit all $R$s for all $n$ equations involved in pairing computation in the first interaction and receive a random challenge $z$. The challenge $z$ can then be used to batch all the $Q$s into a single polynomial as a consecutive powers random linear combination like this,

$$Q_{acc} = \sum_{i=0}^{n - 1}z^i \cdot Q_i(z)$$

The verifier then prepares a new challenge $c$ and checks,

$$\sum_{i=0}^{n - 1}z^i \cdot A_i(c) \cdot B_i(c) \cdots X_i(c) = \sum_{i=0}^{n-1}{z_i \cdot R_i(c) + Q_{acc}(c) \cdot P_{12}(c)}$$

or,

\begin{equation}
    \label{eq:iop-batched}
    \sum_{i=0}^{n - 1}z^i \cdot [A_i(c) \cdot B_i(c) \cdots X_i(c) - R_i(c)] = Q_{acc}(c) \cdot P_{12}(c)
\end{equation}

This reduces $n$ polynomial evaluations of degree ranging from 38 to 76 (as described in Tables \ref{tbl:op:ml-0bit}, \ref{tbl:op:ml-nzbit}, \ref{tbl:op:ml-frob}, \ref{tbl:op:finx}) to a single polynomial of degree 76.

\subsection{Prover}

Prover receives $P$ and $Q$ point pairs, residue witness as computed in \cite{OPP24} and internal input $W$ for 27th root of unity. Prover then computes the $R_i$ and $Q_i$.

\begin{algorithm}
\caption{PolyMulDivQR: Polynomial multiplication and Euclidean division}
\label{alg:prover_polydiv}
\begin{algorithmic}[1]
\Require $P_1, P_2, \ldots, P_n \in \Fq[x]$
\Statex Internal: Divisor $P_{12} \in \Fq[x]$
\Ensure $Q, R \in \Fq[x]$ where $\deg(R) < \deg(P_{12})$ and $\prod_{i=1}^{n} P_i = Q \cdot P_{12} + R$
\State $Q \leftarrow 0$, $R \leftarrow P_1$ \gComm{Initialize quotient and remainder}
\For{$i = 2$ to $n$}
    \State $R \leftarrow R \cdot P_i$ \gComm{Multiply polynomials into R}
\EndFor
\State $d_P \leftarrow \deg(P_{12})$, $d_R \leftarrow \deg(R)$, $l_P \leftarrow \text{leading\_coeff}(P_{12})$
\While{$d_R \geq d_{P}$}
\State $t \leftarrow \text{leading\_coeff}(R)/l_P \cdot x^{d_R - d_{P}}$
\State $Q \leftarrow Q + t$, $R \leftarrow R - t \cdot P_{12}$, $d_R \leftarrow \deg(R)$
\EndWhile
\State \Return $(Q, R)$ \gComm{Quotient and Remainder with $\deg(R) < \deg(P_{12})$}
\end{algorithmic}
\end{algorithm}

\begin{algorithm}
\caption{Pairing Prover: Preparing polynomials}
\label{alg:ppp:2r:prover}
\begin{algorithmic}[1]
\Require $P \in \G_1, Q \in \G_2, c \in \Fe{12}, w \in \{0,1,2\}$
\Statex Internal: $W, W^2 \in \Fe{12}$ \Comment{$W$ is 27th root of unity}
\Ensure Arrays $Q_s, R_s$ of quotients and remainders
\color{gray}
\Statex \small{Every operation involving $T, P \text{ or } Q$ is repeated for $T_i, P_i \text{ and } Q_i$ for multi-pairing.}
\color{black}
\Statex Write $s = 6t + 2 = \sum_{i=0}^{L-1} s_i 2^i$ where $s_i \in \{-1, 0, 1\}$ and $s_L = 1$
\State $c^{-1} \leftarrow 1 / c$ \Comment{Inverse Residue witness}
\State $f \leftarrow c^{-1}$ \Comment{Initialize with inverse residue witness}
\State $T \leftarrow Q$ \gComm{$T_i$ for each pair for multi-pairing}
\Statex
\For{$i = L - 2$ to $0$}
    \State $L_\text{T,T} \leftarrow l_{T,T}(P)$ \gComm{Repeat for $T_i$}
    \State $T \leftarrow [2]T$ \gComm{Repeat for $T_i$}
    \If{$s_i = 0$}
			\State $Q, R \leftarrow {\textit{PolyMulDivQR}}(f, f, L_\text{T,T})$ \Comment{Algorithm \ref{alg:prover_polydiv}}
    \ElsIf{$s_i = 1$}
				\State $Q, R \leftarrow {\textit{PolyMulDivQR}}(f, f, c^{-1}, L_\text{T,T}, l_{T,Q}(P))$ \Comment{Algorithm \ref{alg:prover_polydiv}}
        \State $T \leftarrow T + Q$ \gComm{Repeat for $T_i, Q_i$}
    \ElsIf{$s_i = -1$}
				\State $Q, R \leftarrow {\textit{PolyMulDivQR}}(f, f, c, L_\text{T,T}, l_{T,-Q}(P))$ \Comment{Algorithm \ref{alg:prover_polydiv}}
        \State $T \leftarrow T - Q$ \gComm{Repeat for $T_i, -Q_i$}
    \EndIf
    \State $f \leftarrow R$, $Q_s$.\textit{append}($Q$), $R_s$.\textit{append}($R$)
\EndFor
\Statex
\State $Q_1 \leftarrow \pi_p(Q), Q_2 \leftarrow \pi_p^2(Q)$ \gComm{Repeat for $Q'_i, T_i$}
\State $L_{T,Q_1} \leftarrow l_{T,Q_1}(P)$
\State $T \leftarrow T + Q_1$  \gComm{Repeat for $Q_{1i}, T_i$}
\State $Q, R \leftarrow {\textit{PolyMulDivQR}}(f, L_{T,Q_1}, l_{T,-Q_2}(P))$ \Comment{Algorithm \ref{alg:prover_polydiv}}
\State $f \leftarrow R$, $Q_s$.\textit{append}($Q$), $R_s$.\textit{append}($R$)

\Statex

\If{$w = 1$}
    \State $Q, R \leftarrow {\textit{PolyMulDivQR}}(f, W, c^{-q}, c^{q^2}, c^{-q^3})$ \Comment{Algorithm \ref{alg:prover_polydiv}}
\ElsIf{$w = 2$}
    \State $Q, R \leftarrow {\textit{PolyMulDivQR}}(f, W^2, c^{-q}, c^{q^2}, c^{-q^3})$ \Comment{Algorithm \ref{alg:prover_polydiv}}
\EndIf
\State $Q_s$.\textit{append}($Q$), $R_s$.\textit{append}($R$) \gComm{Final $Q_s$ and $R_s$ accumulation, skip $f$}

\Statex

\State \Return $Q_s, R_s$
\end{algorithmic}
\end{algorithm}

\subsection{Verifier}

Verifier receives all the remainders, $R_i$, and RLC of quotients - $\sum_{i=1}^{n}c_i  \cdot Q_i$.
The verifier then prepares all $c_i$ solely from the remainder coefficients. For the evaluation point $x$, she includes all coefficients from the quotient RLC in the Fiat-Shamir hash as well. This guarantees soundness in Fiat Shamir. Then she uses generated $c_i$ for aggregating $\text{LHS}_i$ evaluated at $x$. In the end, she subtracts the evaluation of quotients RLC at $x$ times $P_{12}(x)$ from the aggregated LHS evaluation. If the result is zero, according to Schwartz--Zippel lemma, all the provided $R_i$ are correct with a very high\footnote{$Pr[P(r_0, r_1, r_2,..., r_d) = 0] \leq \frac{d}{|\F|}$} probability.

\subsection{Performance}

Batching operations in Miller loop steps yielded substantial performance improvements. Processing all operations within a Miller loop step as a single polynomial and performing Schwartz-Zippel verification of polynomial multiplications significantly reduced the overall computational requirements of the pairing verification process.
Combined with residue witness technique for final exponentiation elimination\cite{OPP24}, these optimizations enabled processing of ZK proofs on Starknet even before the modulo builtin\cite{sn_mod_builtin} became available, which later greatly improved finite field arithmetic operations. These techniques were subsequently integrated into Garaga, which remains the most efficient pairing implementation by operation count on Starknet to date.
% https://github.com/shramee/cairo_pairing/blob/main/packages/bn254_u256/src/pairing/schzip/steps.cairo

\subsection{[wip] Protocol Overview}

\subsection{[wip] Single-Round Protocol}

\subsubsection{[wip] Protocol description (Algorithm)}
\subsubsection{[wip] Complexity analysis}


\subsection{Two-Round Protocol with RLC Batching}

\subsubsection{[wip] Protocol description (Algorithm)}
\subsubsection{[wip] Complexity analysis}
\subsubsection{[wip] Comparison with single-round}


\subsection{Three-Round Protocol with Native Field Optimization}

\subsubsection{[wip] Protocol description (Algorithm)}
\subsubsection{[wip] Complexity analysis}
\subsubsection{[wip] Architecture-specific benefits}
